\section{Four marked boundary points and the exploration drawing}\label{sec:pure}

Proposition~\ref{prop:polygon-comparison} compares the random model with
the pure model on the same finite graph. We now identify the pure limit
and specify the finite exploration laws that will be used after stopping.
The distinction between an actual FK crossing and an interface pairing is
important: deterministic boundary closures give the two events different
relative weights.

\subsection{Pairings and boundary closures}

At each lattice edge, the medial representation has a tile with four
ports and two noncrossing pairings of those ports. The open state joins
the two primal corners; the other state joins the dual corners.
Joining the ports of adjacent tiles gives loops and, at marked boundary
changes, open strands. Give an open switch local weight $v_e/\sqrt2$
and a closed switch weight $1$. The \emph{open-cap law} assigns a
configuration weight equal to the product of its local switch weights
times $(\sqrt2)^{\ell}$, where $\ell$ is the number of internal closed
loops. Open strands receive no extra factor. At the pure critical odds
$v_e=\sqrt2$, the two local weights are equal.

Label the four ports in cyclic order by $a,b,c,d$, with designated
primal arcs $ab$ and $cd$. There are two noncrossing pairings:
$E_{\mathrm P}=\{a\leftrightarrow d,\ b\leftrightarrow c\}$ and
$E_{\mathrm D}=\{a\leftrightarrow b,\ c\leftrightarrow d\}$.
Under separate primal closure, $E_{\mathrm P}$ is exactly the event of
an actual primal crossing; under separate dual closure,
$E_{\mathrm D}$ is the analogous dual event.

A four-change exploration also specifies a distinguished starting
port and follows the unique strand issuing from that port until it
reaches another marked port. This choice is part of the experiment,
independent of the cyclic labels used for the electrodes. In the
short-interval experiment of Section~\ref{sec:identification}, the
chosen port is the original chordal starting point.

This is a finite measure defined by the loop representation. It need not
be identified with a boundary partition of the original primal graph.
The following exact change of measure supplies the connection.

\begin{lemma}[Cap change]\label{lem:cap-change}
Let $p$ be the probability of $E_{\mathrm P}$ in a four-change open-cap
law with arbitrary positive FK edge odds. Under separate primal closure,
the probability of an actual primal crossing is
\begin{equation}\label{eq:cap-change}
 \mathcal G(p)=\frac{p}{p+\sqrt2(1-p)}.
\end{equation}
The same formula holds for conditional probabilities after any exploration
history that reveals only encountered switches. The density between the
two full-configuration laws and its reciprocal are bounded by absolute
constants. The analogous statement holds with primal and dual interchanged.
\end{lemma}

\begin{proof}
Fix the internal switches. The two possible boundary pairings acquire
numbers of additional loops differing by one under separate primal
closure. The noncrossing outcome acquires the extra factor $\sqrt2$.
The local edge factors are unchanged. Thus the separate-closure law has
density proportional to
$\ind_{E_{\mathrm P}}+\sqrt2\ind_{E_{\mathrm D}}$
relative to the open-cap law. Normalization proves
\eqref{eq:cap-change} and the density bounds. Conditioning this same
two-valued density on the full revealed history proves the conditional
identity. In particular, no Markov property of a putative limiting curve
is being used. This is the finite cap calculation of~\cite[Section
``Alternating tests after exploration'']{C}; it does not require constant
edge odds.
\end{proof}

\subsection{The pure four-point limit on a flat disk}

For a marked conformal disk $(P;a,b,c,d)$, let $\psi:P\to\Hh$ send
$(a,c,d)$ to $(0,1,\infty)$ and put $\chi=\psi(b)\in(0,1)$.
The primal electrodes are $ab$ and $cd$. With the convention above,
$E_{\mathrm P}$ is rare when $\chi\downarrow0$. Define
\begin{equation}\label{eq:four-point-functions}
 \begin{split}
 f(\chi)&=
 \frac{\sqrt{1-\sqrt{1-\chi}}}
 {\sqrt{1-\sqrt\chi}+\sqrt{1-\sqrt{1-\chi}}},\\
 g(\chi)&=\frac{f(\chi)}{f(\chi)+\sqrt2(1-f(\chi))}.
 \end{split}
\end{equation}
The function $f$ is the pure open-cap pairing probability; $g$ is the
actual crossing probability under separate primal wiring. The
normalization agrees with the $q=2$ specialization in~\cite[Proposition
``Jordan-quad transfer'']{C}. In the convention of
\cite[Equation~(1.1)]{CS}, setting $u=1-\chi$ gives the complementary
pairing probability $1-f(\chi)$. This accounts for the apparent reversal
of the small-interval limit between the two conventions. In particular,
\begin{equation}\label{eq:four-point-properties}
 f(1-\chi)=1-f(\chi),\qquad
 f(\chi)\sim\sqrt{\chi/2}\quad(\chi\downarrow0).
\end{equation}
Both functions are continuous on $[0,1]$ and smooth in its interior.

\begin{proposition}[Pure four-point limit]\label{prop:pure-four-mark}
Let $P$ be a fixed collared orthogonal polygon that is a topological disk
in an oriented flat surface with a locally injective translation projection
to the plane. Put four distinct marks in straight boundary portions.
Approximate $P$ by its lifted square grid at mesh $\delta\downarrow0$,
with critical odds $\sqrt2$ on every edge. Then the open-cap probability
of $E_{\mathrm P}$ tends to $f(\chi_P)$, and the probability of an actual
crossing between separately wired arcs $ab$ and $cd$ tends to $g(\chi_P)$.
The conclusion also holds for compatible medial approximations differing
by an incident boundary layer.
\end{proposition}

\begin{proof}
For a planar polygon this is the four-mark observable theorem of
Chelkak--Smirnov~\cite[Theorem~6.1]{CS}, with boundary pairings and
deterministic closures converted as in Lemma~\ref{lem:cap-change}.
We explain why its proof applies to $P$. This is necessary because the
projection of $P$ need not be injective.

\emph{The finite observable.}
The construction combines the two Dobrushin observables obtained by
joining $c$ to $d$, or $a$ to $d$, with a deterministic exterior arc. Its local
s-holomorphic identities and its relation to the two pairing weights
are identities on tiles. The primitive of its square is single-valued
on the disk and is constant on each of the four boundary arcs. In
cyclic order its normalized boundary values are
$0,1,\kappa_\delta,\kappa_\delta$. If
$p_\delta=\P_{\mathrm{open}}(E_{\mathrm P})$, the exact relation is
\begin{equation}\label{eq:primitive-parameter}
 \kappa_\delta=
 \left(\frac{t_\delta^2+\sqrt2t_\delta}
 {t_\delta^2+\sqrt2t_\delta+1}\right)^2,
 \qquad t_\delta=\frac{1-p_\delta}{p_\delta}.
\end{equation}
This extends continuously to $p_\delta=0,1$ and is one-to-one.
It is the normalization in~\cite[Equations~(6.1)--(6.6)]{CS}, with
their pairing probability equal to $1-p_\delta$.

The phases can be computed in the frame supplied by the projection.
A simple closed curve bounding a disk in this flat surface has total
turning $2\pi$ by Gauss--Bonnet. Closing an observable path by a boundary
arc therefore gives the same phase relations as in the plane.
The collar accommodates the deterministic closing arcs. There is no
additional period for the primitive or ambiguity in the half-differential
on a simply connected polygon.

\emph{Compactness and boundary values.}
The bounds on the primitive use its boundary values, the local boundary
modification by half-rhombi, and the discrete maximum principle.
They do not require a global injective coordinate. Interior estimates
for the observable and its primitive apply in ordinary projection
charts. A finite chart cover of each compact subset consequently gives
subsequential convergence to a bounded harmonic primitive $H$ whose
complex differential has the corresponding holomorphic square root.

The constant data and one-sided normal conditions pass to the limit
near open straight sides by the same boundary estimates. In each such
chart, comparison with distant boundary data first stops at the boundary
of a small chart neighborhood; the comparison function there is bounded.
Its contribution near the side is then suppressed by ordinary planar
discrete harmonic measure. Corners have ordinary wedge charts. Their
isolated values do not affect the bounded harmonic Dirichlet problem.
The marked points lie in straight portions, so their two one-sided
conditions are handled in disjoint neighborhoods in the same way.

\emph{Identification of the remaining constant.}
The sign condition needed here is concrete. Denote the two auxiliary
primitives by $H_\delta^\bullet$ and $H_\delta^\circ$, on the black
and white graphs of the observable construction, respectively. The
black primitive is subharmonic and the white one superharmonic;
these color names follow~\cite{CS}. On the arc $cd$,
$H_\delta^\bullet$ at an inward neighbor is at least its boundary
value $\kappa_\delta$. On $da$, $H_\delta^\circ$ has the reversed inequality.
These follow from the observable's boundary tangent condition and the
half-face modification in~\cite[Section~6]{CS}.

To see that these conditions survive in the required sense, suppose
that a limiting harmonic primitive satisfies $H<\kappa$ on a small
interior contour $\mathcal C$ whose endpoints lie in a straight
subinterval of $cd$. Enclose the region between this contour and the
subinterval, entirely within one chart. Let $V_\delta$ be the discrete
harmonic measure of a smaller subinterval strictly between the endpoints.
It vanishes on the contour. Discrete Green's formula and the inward
inequality give
\begin{equation}\label{eq:local-boundary-sign}
 \sum_{u\in\mathcal C_\delta}
       V_\delta(u_{\rm in})
       \bigl(H_\delta^\bullet(u)-\kappa_\delta\bigr)w_u\ge0,
\end{equation}
where $u_{\rm in}$ is the inward neighbor and $w_u>0$ is the local
conductance; multiple boundary incidences are summed separately.
On a fixed middle portion, $V_\delta(u_{\rm in})\ge c_1\delta$ and
there are order $\delta^{-1}$ incidences. Compact convergence thus
bounds that portion of the sum above by a fixed negative number $-c$.
At distance $r$ from an endpoint, the local harmonic-measure bound is
$V_\delta(u_{\rm in})\le C\delta r^{\eta-1}$ for some $\eta>0$.
Summing over the endpoint portions of length $\mu$, and using the
bounded primitive, gives a contribution at most $C\mu^\eta$.
The remaining compact portion has nonpositive contribution for small
mesh. Thus \eqref{eq:local-boundary-sign}
would imply $0\le-c+C\mu^\eta$, a contradiction for small $\mu$.
The dual argument excludes $H>\kappa$ near a subinterval of $da$.
This is exactly the local comparison in
\cite[Equations~(6.8)--(6.10)]{CS}; all its objects lie in one chart.

Pull the limiting boundary-value problem to a conformal disk.
For $0<\kappa<1$ its solution is the imaginary part of a conformal
map to a horizontal strip of height one with a slit at height $\kappa$,
sending $a$ to the lower left end, $b$ to the right end, and $c$ to the
upper left end. If $d$ mapped to the lower slit bank rather than the
tip, one would have $H<\kappa$ near a part of $cd$; the upper bank
would give $H>\kappa$ near a part of $da$. The preceding contradiction
excludes both. It also excludes $\kappa=0,1$ when the four marks are
distinct. Since the polygon boundary is Jordan in the surface, every
nonempty boundary interval contains a straight subinterval on which
the argument can be applied.

Consequently $d$ maps to the slit tip. In the normalized half-plane
quad $(\Hh;0,\chi_P,1,\infty)$ the strip parameter is
$\kappa=1-\chi_P$. Solving \eqref{eq:primitive-parameter} for the
pairing probability gives \eqref{eq:four-point-functions}.
This identifies every subsequential limit.

\emph{Boundary-layer conventions.}
If the chosen medial graph differs from a primal polygon approximation
by an incident layer, place an easier and a harder orthogonal quad in
the collar, with strict margins. The easier quad extends beyond the
two free sides and ends inside the two wired sides. Restrict the target
and test to their shared grid edges. The target's unmatched pieces
attach only at the two electrodes and cannot join the opposite exterior
zones. Their induced partition is therefore dominated by separate
wiring of the easier electrodes. Every target crossing contains a
shared-edge test crossing. This gives the upper comparison by FK
monotonicity. Interchanging the roles with the harder quad gives the
lower comparison. These are intrinsic disk statements on the same
lifted grid, as in~\cite[Proposition ``Jordan-quad transfer'']{C}.
Let the mesh tend to zero with margins fixed, then shrink the margins.
Marked Jordan convergence in the collared disk gives convergence of
the two moduli to $\chi_P$, completing the proof.
\end{proof}

\begin{corollary}[Fixed quenched tests]\label{cor:fixed-tests}
For each fixed polygon $P$ in Proposition~\ref{prop:pure-four-mark},
assign to every lifted edge the random odds of its projected edge.
At the critical inverse temperature of Theorem~\ref{thm:main}, the
separately wired actual crossing probability tends in environmental
probability to $g(\chi_P)$. The open-cap probability tends to $f(\chi_P)$.
\end{corollary}

\begin{proof}
Combine Proposition~\ref{prop:polygon-comparison} with
Proposition~\ref{prop:pure-four-mark}, then invert the strictly increasing
continuous map $\mathcal G$ in \eqref{eq:cap-change}.
\end{proof}

We fix a countable stock $\mathscr T$ of these tests. A member is
specified by finitely many rational rectangular projection charts,
rational rectilinear cells in those charts, and finite incidence data
recording which copies of overlapping cells belong to the same sheet.
The selected cells form a collared disk, with four marks in straight
portions. Include the dual versions and all finite parity and boundary-layer
placements needed for the medial drawing. There are countably many
finite descriptions of this kind. Every compact collared polygon away
from finitely many punctures can be approximated, with prescribed strict
margins, by such a test. Crucially, a description records the same
projected grid edges on every occurrence of a sheet. It never replaces
the projected disorder by independent copies.

\subsection{A drawing compatible with conditional FK laws}

The usual medial curve can visit a lattice point more than once. For
the slit-domain geometry we separate its local passages by a fixed
deterministic convention. The convention must retain the original
conditional FK law; an arbitrary simple perturbation of a curve would
not suffice.

\begin{lemma}[Compatible drawing]\label{lem:compatible-drawing}
The exploration in a filled lattice domain has a tile-and-collar drawing
in a Jordan polygon $V_n$ with the following properties.
\begin{enumerate}[(i)]
\item The marked boundary of $V_n$ has the same uniform Jordan limit
as that of $D_n$.
\item The drawn strand $\widehat\gamma_n$ is simple, follows the same
ordered local pieces as $\gamma_n$, and satisfies
$d_{\mathrm{curv}}(\widehat\gamma_n,\gamma_n)=O(\delta_n)$.
\item After a switch-revealing stop, the remaining domain is a simple
slit disk. Under deterministic cap closure the unknown edges have
their original FK odds, with the revealed open edges contracted and
revealed closed edges deleted.
\item Each bank of the prefix has physical access, within $O(\delta_n)$,
to its corresponding known edges and boundary wire. Every completion
embeds its actual primal and dual paths without crossing the past.
No uncertain physical edge is duplicated.
\end{enumerate}
The same drawing and caps can be used away from an interval at which
two extra boundary changes are inserted.
\end{lemma}

\begin{proof}
We recall the finite construction of~\cite[Lemma ``A compatible simple
drawing'']{C}. Make one tile for every edge, including boundary edges.
Glue the four triangles directed toward an interior square center.
For each boundary edge, keep its exterior triangle as a separate tab
with its own exterior dual vertex. Different exterior tabs are not
glued along their radial sides; primal boundary vertices retain their
interior incidences. Flatten the tabs and add a thinner Jordan collar.
The displacement is $O(\delta_n)$, and tracing the tabs in the original
boundary order proves (i).

Draw the two possible switches as disjoint local arcs. Draw boundary
caps and the wires behind them in the collar; unmatched ports lead to
the marked outer points. There are finitely many local choices, all
made deterministically. The switches follow the prescribed medial
pairings. Where a wired boundary edge is bypassed by a cap, either
switch state gives the same continuation of the inside ports, up to
a microscopic detour and possibly a microscopic loop. The edge still
retains its random state in the FK law. On the free side, the cap
around a primal vertex gives the turn beside the dual boundary wire.
Match the ordered pieces of the original and drawn strands. Each
matched pair lies in a bounded number of neighboring tiles; the
resulting increasing parametrizations prove (ii).

The tile Euler identity gives the FK component weight under each
deterministic closure. Revealing an encountered switch therefore does
exactly the usual contraction or deletion and leaves the other edge
odds unchanged. This proves (iii), including a stop inside a tile if
that entire encountered tile is revealed.

On an open primal turn the bank uses the actual revealed edge; on the
other turn it stays at a primal corner. Along a cap it stays at a
vertex or follows vertices of the same wire interval. These pieces
give the asserted short physical accesses. The dual statement is
identical. The embedding of every completion follows from the local
noncrossing pairings. Projection back to the original grid moves
points by $O(\delta_n)$ and duplicates no edge, although separate
boundary copies of a vertex can project to one physical vertex.
This proves (iv). None of the construction uses equality of edge odds.
\end{proof}

All later stopping rules reveal only encountered switches. If a stop
occurs during a drawn local piece, its tile is already known; completing
that tile changes the geometric prefix only within $O(\delta_n)$.
The filtration can therefore include capacity times as well as
switch times without querying a future switch.
